\documentclass[12pt,letterpaper]{article}
\usepackage{graphicx,textcomp}
\usepackage{natbib}
\usepackage{setspace}
\usepackage{fullpage}
\usepackage{color}
\usepackage[reqno]{amsmath}
\usepackage{amsthm}
\usepackage{fancyvrb}
\usepackage{amssymb,enumerate}
\usepackage[all]{xy}
\usepackage{endnotes}
\usepackage{lscape}
\newtheorem{com}{Comment}
\usepackage{float}
\usepackage{hyperref}
\newtheorem{lem} {Lemma}
\newtheorem{prop}{Proposition}
\newtheorem{thm}{Theorem}
\newtheorem{defn}{Definition}
\newtheorem{cor}{Corollary}
\newtheorem{obs}{Observation}
\usepackage[compact]{titlesec}
\usepackage{dcolumn}
\usepackage{tikz}
\usetikzlibrary{arrows}
\usepackage{multirow}
\usepackage{xcolor}
\newcolumntype{.}{D{.}{.}{-1}}
\newcolumntype{d}[1]{D{.}{.}{#1}}
\definecolor{light-gray}{gray}{0.65}
\usepackage{url}
\usepackage{listings}
\usepackage{color}

\definecolor{codegreen}{rgb}{0,0.6,0}
\definecolor{codegray}{rgb}{0.5,0.5,0.5}
\definecolor{codepurple}{rgb}{0.58,0,0.82}
\definecolor{backcolour}{rgb}{0.95,0.95,0.92}

\lstdefinestyle{mystyle}{
	backgroundcolor=\color{backcolour},   
	commentstyle=\color{codegreen},
	keywordstyle=\color{magenta},
	numberstyle=\tiny\color{codegray},
	stringstyle=\color{codepurple},
	basicstyle=\footnotesize,
	breakatwhitespace=false,         
	breaklines=true,                 
	captionpos=b,                    
	keepspaces=true,                 
	numbers=left,                    
	numbersep=5pt,                  
	showspaces=false,                
	showstringspaces=false,
	showtabs=false,                  
	tabsize=2
}
\lstset{style=mystyle}
\newcommand{\Sref}[1]{Section~\ref{#1}}
\newtheorem{hyp}{Hypothesis}

\title{Problem Set 2}
\date{Due: February 18, 2026}
\author{Applied Stats II}


\begin{document}
	\maketitle
	\section*{Instructions}
	\begin{itemize}
		\item Please show your work! You may lose points by simply writing in the answer. If the problem requires you to execute commands in \texttt{R}, please include the code you used to get your answers. Please also include the \texttt{.R} file that contains your code. If you are not sure if work needs to be shown for a particular problem, please ask.
		\item Your homework should be submitted electronically on GitHub in \texttt{.pdf} form.
		\item This problem set is due before 23:59 on Wednesday February 18, 2026. No late assignments will be accepted.

	\end{itemize}

	\vspace{.25cm}

\noindent We're interested in what types of international environmental agreements or policies people support (\href{https://www.pnas.org/content/110/34/13763}{Bechtel and Scheve 2013)}. So, we asked 8,500 individuals whether they support a given policy, and for each participant, we vary the (1) number of countries that participate in the international agreement and (2) sanctions for not following the agreement. \\

\noindent Load in the data labeled \texttt{climateSupport.RData} on GitHub, which contains an observational study of 8,500 observations.

\begin{itemize}
	\item
	Response variable: 
	\begin{itemize}
		\item \texttt{choice}: 1 if the individual agreed with the policy; 0 if the individual did not support the policy
	\end{itemize}
	\item
	Explanatory variables: 
	\begin{itemize}
		\item
		\texttt{countries}: Number of participating countries [20 of 192; 80 of 192; 160 of 192]
		\item
		\texttt{sanctions}: Sanctions for missing emission reduction targets [None, 5\%, 15\%, and 20\% of the monthly household costs given 2\% GDP growth]
		
	\end{itemize}
	
\end{itemize}

\newpage
\noindent Please answer the following questions:

\begin{enumerate}
	\item
	Remember, we are interested in predicting the likelihood of an individual supporting a policy based on the number of countries participating and the possible sanctions for non-compliance.
	\begin{enumerate}
		\item [] Fit an additive model. Provide the summary output, the global null hypothesis, and $p$-value. Please describe the results and provide a conclusion.
			\begin{description}
			\item Null hypothesis: The number of countries and sanctions for missing emission reduction targets have no effect on the policy choice.
		\end{description}
		\begin{description}
			\item Alternative hypothesis: There is at least one variable that effects the policy choice. 
		\end{description}
	\lstinputlisting[language=R, firstline=50,lastline=53]{PS02_stats.R} 
\begin{table}[!htbp] \centering 
	\caption{Additive Model} 
	\label{} 
	\begin{tabular}{@{\extracolsep{5pt}}lc} 
		\\[-1.8ex]\hline 
		\hline \\[-1.8ex] 
		& \multicolumn{1}{c}{\textit{Dependent variable:}} \\ 
		\cline{2-2} 
		\\[-1.8ex] & choice \\ 
		& Coefficients \\ 
		\hline \\[-1.8ex] 
		countries80 of 192 & 0.336$^{***}$ \\ 
		& (0.054) \\ 
		& \\ 
		countries160 of 192 & 0.648$^{***}$ \\ 
		& (0.054) \\ 
		& \\ 
		sanctions5\% & 0.192$^{***}$ \\ 
		& (0.062) \\ 
		& \\ 
		sanctions15\% & $-$0.133$^{**}$ \\ 
		& (0.062) \\ 
		& \\ 
		sanctions20\% & $-$0.304$^{***}$ \\ 
		& (0.062) \\ 
		& \\ 
		Constant & $-$0.273$^{***}$ \\ 
		& (0.054) \\ 
		& \\ 
		\hline \\[-1.8ex] 
		Observations & 8,500 \\ 
		Log Likelihood & $-$5,784.130 \\ 
		Akaike Inf. Crit. & 11,580.260 \\ 
		\hline 
		\hline \\[-1.8ex] 
		\textit{Note:}  & \multicolumn{1}{r}{$^{*}$p$<$0.1; $^{**}$p$<$0.05; $^{***}$p$<$0.01} \\ 
	\end{tabular} 
\end{table} 
\newpage
	\lstinputlisting[language=R, firstline=61,lastline=63]{PS02_stats.R} 
\begin{table}[!htbp] \centering 
	\caption{Null Model} 
	\label{} 
	\begin{tabular}{@{\extracolsep{5pt}}lc} 
		\\[-1.8ex]\hline 
		\hline \\[-1.8ex] 
		& \multicolumn{1}{c}{\textit{Dependent variable:}} \\ 
		\cline{2-2} 
		\\[-1.8ex] & choice \\ 
		& Coefficients \\ 
		\hline \\[-1.8ex] 
		Constant & $-$0.007 \\ 
		& (0.022) \\ 
		& \\ 
		\hline \\[-1.8ex] 
		Observations & 8,500 \\ 
		Log Likelihood & $-$5,891.705 \\ 
		Akaike Inf. Crit. & 11,785.410 \\ 
		\hline 
		\hline \\[-1.8ex] 
		\textit{Note:}  & \multicolumn{1}{r}{$^{*}$p$<$0.1; $^{**}$p$<$0.05; $^{***}$p$<$0.01} \\ 
	\end{tabular} 
\end{table} 
\lstinputlisting[language=R, firstline=71,lastline=77]{PS02_stats.R} 
		\begin{description}
			\item The anova test above shows a p value lower that 0.05, therefore, we reject the null hypothesis that countries and sanctions have zero effect on the policy choice. At least one of the predictors is significant in predicting the policy choice.  
		\end{description}
	\end{enumerate}
	\item
	If any of the explanatory variables are statistically significant in this model, then:
	\begin{enumerate}
		\item
		For the policy in which nearly all countries participate [160 of 192], how does increasing sanctions from 5\% to 15\% change the odds that an individual will support the policy? (Interpretation of a coefficient)
		\begin{equation}
	-0.133 - 0.192 = -0.325
\end{equation}
\begin{description}
	\item When 160 of 192 countries participate, increasing sanctions from 5 \% to 15 \% reduces the log odds of supporting the policy by 0.325, holding all elseconstant. 
\end{description}
		\item
		For the policy in which very few countries participate [20 of 192], how does increasing sanctions from 5\% to 15\% change the odds that an individual will support the policy? (Interpretation of a coefficient)
		\begin{equation}
			-0.133 - 0.192 = -0.325
		\end{equation}
			\begin{description}
			\item When 20 of 192 countries participate, increasing sanctions from 5 \% to 15 \% reduces the log odds of supporting the policy by 0.325. Which is the same as the answer for part a because there is no interaction in this additive model between number of countries participating and percent of sanctions. 
		\end{description}
		\item
		What is the estimated probability that an individual will support a policy if there are 80 of 192 countries participating with no sanctions? 
	\begin{description}
		\item[Calculating the log odds given 80 of 192 countries and no sanctions ] 
	\end{description}
	\begin{equation}
	-0.273 + 0.336 = 0.063
\end{equation}
	\begin{description}
	\item[Calculating the probability] 
\end{description}
	\lstinputlisting[language=R, firstline=80, lastline=82]{PS02_stats.R} 
\begin{description}
	\item If there are 80 of 192 countries participating and no sanctions, the probability of supporting the policy is about 51.6 \%. 
\end{description}
	\end{enumerate}
	\item
	Would the answers to 2a and 2b potentially change if we included an interaction term in this model? Why? 
	\begin{itemize}
		\item Perform a test to see if including an interaction is appropriate.
	\end{itemize}
	\begin{description}
		\item Null hypothesis: there is no interaction effect between the number of countries and sanctions for missing emission reduction targets on the choice on the policy. 
	\end{description}
		\begin{description}
		\item Alternative hypothesis: There is at least one interaction effect between countries and sanctions. 
	\end{description}
		\lstinputlisting[language=R, firstline=85,lastline=89]{PS02_stats.R} 
	\begin{table}[!htbp] \centering 
		\caption{Interaction Model} 
		\label{} 
		\begin{tabular}{@{\extracolsep{5pt}}lc} 
			\\[-1.8ex]\hline 
			\hline \\[-1.8ex] 
			& \multicolumn{1}{c}{\textit{Dependent variable:}} \\ 
			\cline{2-2} 
			\\[-1.8ex] & choice \\ 
			& Coefficients \\ 
			\hline \\[-1.8ex] 
			countries80 of 192 & 0.376$^{***}$ \\ 
			& (0.106) \\ 
			& \\ 
			countries160 of 192 & 0.613$^{***}$ \\ 
			& (0.108) \\ 
			& \\ 
			sanctions5\% & 0.122 \\ 
			& (0.105) \\ 
			& \\ 
			sanctions15\% & $-$0.097 \\ 
			& (0.108) \\ 
			& \\ 
			sanctions20\% & $-$0.253$^{**}$ \\ 
			& (0.108) \\ 
			& \\ 
			countries80 of 192:sanctions5\% & 0.095 \\ 
			& (0.152) \\ 
			& \\ 
			countries160 of 192:sanctions5\% & 0.130 \\ 
			& (0.151) \\ 
			& \\ 
			countries80 of 192:sanctions15\% & $-$0.052 \\ 
			& (0.152) \\ 
			& \\ 
			countries160 of 192:sanctions15\% & $-$0.052 \\ 
			& (0.153) \\ 
			& \\ 
			countries80 of 192:sanctions20\% & $-$0.197 \\ 
			& (0.151) \\ 
			& \\ 
			countries160 of 192:sanctions20\% & 0.057 \\ 
			& (0.154) \\ 
			& \\ 
			Constant & $-$0.275$^{***}$ \\ 
			& (0.075) \\ 
			& \\ 
			\hline \\[-1.8ex] 
			Observations & 8,500 \\ 
			Log Likelihood & $-$5,780.983 \\ 
			Akaike Inf. Crit. & 11,585.970 \\ 
			\hline 
			\hline \\[-1.8ex] 
			\textit{Note:}  & \multicolumn{1}{r}{$^{*}$p$<$0.1; $^{**}$p$<$0.05; $^{***}$p$<$0.01} \\ 
		\end{tabular} 
	\end{table} 
	\newpage
	\lstinputlisting[language=R, firstline=98,lastline=104]{PS02_stats.R} 
	\begin{description}
		\item Based on the anova test above, the result (0.3912) is larger than the p value of 0.05, therefore, we fail to reject the null hypothesis that there is no interaction between countries and sanctions. Thus, the additive model is better at explaining the relationship between the variables and the effect of sanctions is constant across the amount of countries participating. 
	\end{description}
	\end{enumerate}


\end{document}
